\section{Why There Is No Preferred Frame}

This chapter is about a problem that kills most theories like this one, and how the curved crystal survives it. The problem is Lorentz invariance, the deep fact that the laws of physics look the same no matter how fast you are moving, with no special ``at rest'' frame anywhere. It sounds technical. It is actually the make-or-break test for any theory that builds space from a discrete structure.

\subsection{The grid problem}

Suppose you tried to build space from an ordinary flat grid, like graph paper, cells in neat rows and columns. It seems natural. But it has a fatal flaw: a grid has preferred directions. Along the rows and columns is different from along the diagonals. On such a grid, the speed of light would depend on which way it traveled, faster along the rows, slower on the diagonal, say. And it would depend on energy, with high-energy light feeling the grain of the grid.

Experiment forbids this, to staggering precision. Light from distant gamma-ray bursts, after traveling billions of years, arrives with high-energy and low-energy photons still neck and neck. If space had a grid grain, they would have drifted apart. They have not. So any theory whose space has built-in preferred directions is simply ruled out. This is where most discrete-space ideas die.

\subsection{The surprising rescue: curvature}

The crystal escapes, and the reason is beautiful and not obvious. It is not random. It is an exact, ordered crystal, every cell identical. You would think an ordered crystal is the worst case, the most grid-like of all. The opposite is true, because the crystal lives in curved hyperbolic space.

Here is the heart of it. In flat space, a direction can be carried unchanged across the whole space. ``North'' here is the same as ``north'' a mile away. So a flat grid's preferred directions are globally defined, real, and detectable. But in curved space, you cannot carry a direction unchanged from one place to a distant place. Curvature scrambles directions as you move. So even though the crystal is perfectly regular, its directions do not line up across the space. From any local spot, the crystal's links point every which way, indistinguishable from a random foam with no preferred direction at all.

The slogan is: curvature, not disorder, is what removes the preferred frame. You do not need a messy random space to hide the grain. You need a curved one. The crystal can be an exact, beautiful, ordered honeycomb and still look, to anyone living inside it, perfectly smooth and directionless.

\subsection{The test, with numbers}

The simulation checks this directly. Take a flat lattice and measure how much the speed of a signal depends on direction: the dependence is large, and it grows with energy, exactly the fatal grain. Now measure the same thing on the hyperbolic crystal: the direction-dependence is consistent with zero, and it shrinks as the crystal grows. The crystal predicts no first-order grain at all.

That turns into a real, passable test. The gamma-ray-burst timing that rules out a flat lattice is comfortably passed by the curved crystal, because the curved crystal predicts the effect simply is not there. The same observation that kills the grid leaves the crystal standing. And it cuts the other way too: this is a genuine prediction. If anyone ever did detect a first-order direction- or energy-dependence of light's speed, the curved-crystal picture would be falsified.

\subsection{Why this chapter matters}

It would be easy to wave at ``space emerges from a web'' and never face this. Most attempts cannot face it. The grain gives them away. The curved crystal faces it and passes, and it passes precisely because of the hyperbolic geometry we spent Part II building. The curvature was not chosen for beauty or for room alone. It is also exactly what buys Lorentz invariance, the symmetry our universe insists on. Three different needs, room to grow, the right symmetries, and now Lorentz safety, are all met by the single choice of negative curvature.

\textbf{Takeaway:} a flat grid has preferred directions that experiment rules out, which kills most discrete-space theories. The crystal survives because curvature scrambles directions, so an exact ordered hyperbolic crystal looks directionless from inside. Curvature, not disorder, buys Lorentz invariance, and it yields a real, falsifiable prediction of no grain in the speed of light.
